\chapter{Adaptive Streaming}

\textbf{Course:} Multimedia Communications, Universita di Padova, Dipartimento di Ingegneria dell'Informazione.

\textbf{Date:} 20/05/2026.

\bigskip\hrule\bigskip

\section{Lecture Overview}

This lecture covers four main topics:

\begin{enumerate}
\item \textbf{Protocol evolution:} the technical transition from traditional network-driven \textit{push} architectures, such as RTP over UDP, to modern client-driven \textit{pull} adaptive infrastructures running over HTTP.
\item \textbf{Modern transport optimization:} the impact of QUIC and HTTP/3 on high-bitrate multimedia streaming.
\item \textbf{Analytical system modeling:} the continuous-time mathematical fluid model used to describe playout buffer behavior.
\item \textbf{User experience engineering:} Quality of Experience (QoE) metrics and design trade-offs inside \textit{Adaptive Bitrate} (ABR) decision loops.
\end{enumerate}

\section{Classification of Multimedia Services}

Multimedia traffic can be classified by its technical requirements and tolerance to network impairments.

\begin{center}
\begin{tabular}{|l|l|l|l|}
\hline
\textbf{Service category} & \textbf{Real-world examples} & \textbf{Critical requirement} & \textbf{Delay tolerance} \\ \hline
Bulk transfer & Social media photos, messaging & Data integrity, no loss & High, seconds \\ \hline
Video on demand & Netflix, YouTube, Disney+ & Continuity, no jitter & Medium, 1-3 s startup \\ \hline
Live streaming & Twitch, DAZN, YouTube Live & Scalability and recency & Low, below 10 s \\ \hline
Real-time interactive & Teams, Zoom, WhatsApp calls & Ultra-low latency & Critical, below 150 ms \\ \hline
Cloud gaming / VR & GeForce Now, Meta Quest & Interaction and bitrate & Extreme, below 50 ms \\ \hline
\end{tabular}
\end{center}

\section{Video Communication Through Networks}

An uncompressed video source emits bits at a \textbf{constant rate}. For example, uncompressed YCbCr HD video with $1920 \times 1080$ pixels, frame rate $50$ Hz, and $12$ bits per pixel for 4:2:0 sampling requires about $1.244$ Gbps.

After compression, this rate is much lower. Typical rates for this example are in the range $5 \div 20$ Mbps. However, after compression the coding rate is \textbf{no longer constant}: I frames are much heavier than P and B frames.

The compressed bitrate becomes approximately stable only when averaged at the \textbf{Group of Pictures (GOP)} level.

\begin{center}\includegraphics[width=0.806\linewidth]{Pasted_image_20260520131140.png}\end{center}
\textit{Figure 1: Scatter plot of frame size in bits versus frame index for a rate-controlled H.264 encoder, showing I frames around $10^5$ bits, P frames around $10^4$ bits, and B frames around $10^4$ bits on a logarithmic scale.}

\section{Video Content Distribution Challenges}

The distribution paradigm is a \textbf{one-to-many architecture}:

\begin{itemize}
\item \textbf{Single source:} a central video server hosts and transmits the master content.
\item \textbf{Mass destinations:} millions of concurrent, independent end users receive the content.
\item \textbf{User heterogeneity:} each user differs in network path, temporal channel dynamics, hardware, and display format.
\end{itemize}

The four dimensions of user heterogeneity are:

\begin{center}
\begin{tabular}{|l|l|}
\hline
\textbf{Dimension} & \textbf{Description} \\ \hline
Network paths & Different link capacities, asymmetric propagation delays, and routing behaviors. \\ \hline
Temporal dynamics & Time-varying channel conditions, dominated by fast fading and shadowing in wireless/mobile links. \\ \hline
Hardware constraints & Different local memory sizes, battery lifetimes, and hardware decoding capabilities. \\ \hline
Display formats & Devices range from low-resolution smartphone screens to 4K/8K TV sets. \\ \hline
\end{tabular}
\end{center}

\begin{center}\includegraphics[width=0.806\linewidth]{Pasted_image_20260520131216.png}\end{center}
\textit{Figure 2: One-to-many distribution diagram with one video server connected to a 4K smart TV over stable FTTH fiber, a laptop over variable shared Wi-Fi, and a smartphone over low/fading 5G with severe battery limits.}

\section{The Conversational Paradigm}

Conversational services include:

\begin{itemize}
\item \textbf{One-to-one:} video calls. The challenge is that channel capacity fluctuates continuously, so the encoder must adapt the bitrate live with almost no buffering delay.
\item \textbf{Few-to-few:} group conferencing. Calls are routed through a central \textit{Selective Forwarding Unit} (SFU) or \textit{Multipoint Control Unit} (MCU).
\end{itemize}

The live engineering bottleneck is the same heterogeneity seen in streaming, but with a stricter constraint: the system cannot rely on pre-encoded files. It must encode and distribute live video to users with different capacities under extreme \textbf{sub-150 ms} latency constraints.

\begin{center}\includegraphics[width=0.806\linewidth]{Pasted_image_20260520131236.png}\end{center}
\textit{Figure 3: Conversational diagram with Speaker A on fiber and a 4K camera sending to a conference server, which forwards to Listener B over Wi-Fi and Listener C over 4G.}

\section{The Coding Rate Bottleneck}

\begin{calloutbox}{red!70!black}{Coding Rate Constraint}
\textbf{Statement:} To ensure continuous playout or fluid interaction, the video coding bitrate $R_C$ must satisfy the channel capacity limit $C_k$ for user $k$:

\[
R_C \leq C_k
\]

\textbf{Interpretation:} if the selected coding rate is larger than the available capacity, packets accumulate, delay grows, losses occur, and playback can freeze.
\end{calloutbox}

Because each client sees a different $C_k$, a single static $R_C$ creates a symmetric dilemma:

\begin{itemize}
\item \textbf{Cliff effect:} if $R_C > C_k$, packets accumulate in network buffers, causing delay, drops, and severe playback freezes.
\item \textbf{Leveling effect:} if $R_C = \min_k C_k$, high-capacity users receive unnecessarily low quality, wasting their channel potential.
\end{itemize}

The static delivery model must therefore be replaced by an \textbf{adaptive coding rate} framework.

\section{Compression Fragility and Network Friendliness}

Highly compressed video relies on strict spatial/temporal prediction and CABAC contexts. Standard IP routers drop packets blindly during congestion, which can cause severe spatial and temporal error propagation.

In managed domains, such as \textbf{CDN edge nodes} or \textbf{5G mobile gateways}, specialized middleboxes called \textbf{Media Aware Network Elements (MANEs)} can inspect NAL headers. Modern codecs from H.264 onward separate compression data, the \textit{Video Coding Layer} (VCL), from transport-oriented syntax, the \textit{Network Abstraction Layer} (NAL).

MANEs can perform \textbf{smart dropping}: during congestion, they discard non-reference B frames while protecting critical data such as SPS, PPS, and IDR frames, without decoding the video.

\begin{center}\includegraphics[width=0.968\linewidth]{Pasted_image_20260520131321.png}\end{center}
\textit{Figure 4: Encoder-to-decoder pipeline where NAL units pass through a Media Aware Network Element; the MANE filters less important slice units so an I-B-P-B sequence can be simplified to I-P under congestion.}

\section{Modern Multimedia Stacks Push vs Pull}

Two delivery paradigms are used:

\begin{center}
\begin{tabular}{|l|l|l|l|l|}
\hline
\textbf{Paradigm} & \textbf{Objective} & \textbf{Control model} & \textbf{Dominant stack} & \textbf{Typical use} \\ \hline
Real-time push & Minimize absolute interactive latency, below 150 ms & Server-driven; content is pushed immediately & WebRTC and RTP over UDP & Conversational audio/video \\ \hline
Adaptive streaming pull & Maximize visual QoE and scale to global concurrency & Client-driven; receiver requests blocks based on local status & HTTP/3 over QUIC & VoD and live broadcast \\ \hline
\end{tabular}
\end{center}

\begin{center}\includegraphics[width=0.839\linewidth]{Pasted_image_20260520133139.png}\end{center}
\textit{Figure 5: Stack comparison between real-time push with IP, UDP, RTP, RTCP, WebRTC, and application A/V codec, and stored/live streaming pull with IP, UDP, QUIC, HTTP/3, DASH/HLS, CDN edge caching, and application A/V codec.}

\section{The Real-Time Stack WebRTC RTP and UDP}

\begin{itemize}
\item \textbf{UDP} is a bare-bones protocol that provides multiplexing/demultiplexing through sockets. It avoids connection handshakes and retransmission delays, trading perfect data integrity for minimum latency.
\item \textbf{RTP}, the \textit{Real-time Transport Protocol}, adds sequence numbers to detect lost or out-of-order packets and timestamps for audio/video synchronization and jitter management.
\item \textbf{RTCP}, the \textit{RTP Control Protocol}, periodically sends QoS metrics such as packet loss rate and RTT back to the sender, enabling dynamic $R_C$ adaptation.
\item \textbf{WebRTC}, \textit{Web Real-Time Communication}, wraps RTP/RTCP into a browser-native framework, enforcing mandatory encryption and solving NAT/firewall connectivity problems.
\end{itemize}

Typical use cases are video calls, cloud gaming, and remote control.
\begin{center}\includegraphics[width=0.484\linewidth]{Pasted_image_20260520133201.png}\end{center}
\textit{Figure 6: Real-time push stack with IP, UDP, RTP, RTCP, WebRTC, and application A/V codec.}

\section{One-to-Many Real-Time Communications The SVC Paradigm}

In a video call, the encoder must provide different coding rates for different destinations. \textbf{Scalable Video Coding} (SVC) follows the principle \textit{encode once, decode many}.

SVC compresses the source video into a single, unified bitstream layered by importance:

\begin{itemize}
\item \textbf{Base Layer (BL):} essential low-bitrate data with rate $R_B$. It is highly protected and guarantees baseline quality:
\end{itemize}

\[
R_B \leq \min_k C(k)
\]

\begin{itemize}
\item \textbf{Enhancement Layers (EL):} progressive refinement layers containing detail not present in the base layer.
\end{itemize}

Each layer behaves as an independent network payload sub-stream. Enhancement layers improve fidelity, spatial resolution, or frame rate, but are useless without the base layer.

\begin{center}\includegraphics[width=0.839\linewidth]{Pasted_image_20260520133226.png}\end{center}
\textit{Figure 7: SVC hierarchy with a base layer producing baseline quality, enhancement layer 1 producing medium fidelity, and enhancement layer 2 producing high fidelity, each dependent on previous layers.}

\section{Types of Scalability}

SVC can provide:

\begin{itemize}
\item \textbf{Temporal scalability:} obtained by choosing a temporal prediction structure that supports multiple frame-rate levels.
\item \textbf{Quality scalability:} a quality version of the video is encoded in the base layer and used as prediction for enhancement layers, or progressive coding such as bitplanes is used.
\item \textbf{Spatial scalability:} similar to quality scalability, but the base layer is low-resolution, or multi-resolution transforms are used.
\end{itemize}

A layer can combine multiple scalability types. This is signaled in the NAL header through layer identifiers.
\begin{center}\includegraphics[width=0.839\linewidth]{Pasted_image_20260520133327.png}\end{center}

\textit{Figure 8: Temporal scalability diagram with I frames in the base layer at 5 fps, P frames in enhancement layer 1 at 15 fps, and B frames in enhancement layer 2 at 30 fps.}

\begin{center}\includegraphics[width=0.806\linewidth]{Pasted_image_20260520133342.png}\end{center}
\textit{Figure 9: Quality/spatial scalability block diagram where the input video is encoded into a base layer with quantizer $Q_B$, decoded, subtracted from the input to form a residual, then encoded into an enhancement layer with quantizer $Q_E$.}

\section{SVC Infrastructure and Selective Forwarding}

In multipoint conferencing, the speaker sends the SVC-coded video to a server called a \textbf{Selective Forwarding Unit} (SFU), which acts as a MANE. It inspects NAL header metadata, such as Layer ID and Temporal ID, without touching VCL payload.

Smart layer pruning works as follows:

\begin{itemize}
\item \textbf{High-capacity path:} the SFU forwards the complete layered bitstream, BL + EL1 + EL2.
\item \textbf{Low-capacity path:} the SFU avoids random packet dropping and performs intelligent NAL-level dropping, stripping enhancement layers and forwarding only the base layer.
\end{itemize}

According to the number of layers, adaptation can be more or less fine-grained. The result is adaptation with no server re-encoding overhead and no error propagation.

\begin{center}\includegraphics[width=0.968\linewidth]{Pasted_image_20260520134145.png}\end{center}
\textit{Figure 10: Speaker encodes one SVC stream with BL, EL1, and EL2, sends it to an SFU/MANE, which forwards all layers to a fiber user and only the base layer to a 4G mobile user.}

\section{Why SVC Cannot Be Used for Web-Scale Streaming}

SVC is not suitable for web-scale streaming because of infrastructure, bitrate, caching, and hardware costs:

\begin{itemize}
\item In one-to-few real-time video, one MANE/SFU is computationally viable.
\item At web scale, millions of users would require thousands of interconnected MANEs/routers across the network topology to filter layers per user, destroying economic viability and the stateless routing principle of the Internet.
\item SVC introduces heavy syntactic and signaling overhead, typically 10\% to 25\% higher bitrate than single-layer streams at the same visual quality, due to inter-layer prediction data.
\item SVC uses hierarchically dependent layers. CDNs cannot map these layers efficiently to static, atomic URI cache objects. If a base layer segment is missing, all associated enhancement layers become undecodable.
\item Layered decoding increases computational complexity. Mass-market hardware decoders in smart TVs and low-cost mobile SoCs usually lack native ASIC support for SVC.
\end{itemize}

\section{Why HTTP Works Well for Video Streaming}

HTTP works well for video streaming for three main reasons.

\textbf{Firewall and NAT traversal:} legacy streaming protocols such as RTSP/RTP over UDP use random dynamic ports that are often blocked by home and enterprise firewalls. HTTP uses standard web ports, \textbf{80} and \textbf{443}, so it traverses networks almost universally.

\textbf{Stateless server scalability:} the origin server does not maintain playout clocks, session timers, or client-specific state. It behaves as an optimized stateless file dispenser, reducing CPU and memory overhead.

\textbf{Native caching integration:} slicing video into standard HTTP objects integrates with proxies, edge caches, and \textbf{Content Delivery Networks} (CDNs). Media files are replicated without specialized video-aware software.

The drawback is that classic web traffic runs over \textbf{TCP}, with its rigid state machine.
\begin{center}
\includegraphics[width=1.000\linewidth]{Pasted_image_20260520134208.png}\end{center}
\textit{Figure 11: HTTP streaming path where a video client requests \texttt{chunk\_4.mp4} from a CDN edge node; the edge forwards cache misses to the origin server, stores the chunk, then returns a 200 OK payload to the client.}

\section{The Latency Paradox and TCP Limitations}

HTTP streaming over classic networks relies on \textbf{TCP}, the \textit{Transmission Control Protocol}. TCP guarantees absolute reliability, but this creates a penalty for multimedia delivery.

\textbf{Head-of-line (HOL) blocking:} TCP treats data as one sequential byte stream. If a middle packet, such as Packet 2, is lost, the operating system kernel blocks delivery of subsequent packets even if Packets 3 and 4 have already arrived.

The application layer receives no data while TCP waits for retransmission and repair. This transport freeze can cause buffer depletion and application stalls.

\begin{center}\includegraphics[width=0.806\linewidth]{Pasted_image_20260520134237.png}\end{center}
\textit{Figure 12: TCP kernel buffer with Packet 1 delivered, Packet 2 lost and blocked, and Packets 3 and 4 held in kernel memory although they arrived correctly.}

\section{The Generational Leap QUIC and HTTP3}

HTTP/3 abandons TCP and uses \textbf{UDP} as its foundation. Reliability and congestion control are moved from the operating system kernel to user space through \textbf{QUIC}.

QUIC introduces:

\begin{itemize}
\item \textbf{True stream isolation:} different data components are independent streams. If Stream 2 loses a packet, only Stream 2 pauses, while audio or subtitles can continue without head-of-line blocking.
\item \textbf{0-RTT connection handshake:} transport and cryptographic TLS 1.3 handshakes are combined to reduce startup delay.
\item \textbf{Connection migration:} connections are tracked through unique Connection IDs rather than IP/port pairs, so playback can survive Wi-Fi to 5G handovers without disconnecting.
\end{itemize}

\begin{center}\includegraphics[width=0.806\linewidth]{Pasted_image_20260520134259.png}\end{center}
\textit{Figure 13: One QUIC connection over UDP multiplexing three independent streams: audio, video, and subtitles/manifest updates; loss in the video stream pauses only that stream.}

\section{Content Delivery Networks CDNs}

Centralized content deployment creates high propagation delay, proportional to distance, and core-network bottlenecks.

A \textbf{Content Delivery Network} (CDN) is a global grid of geographically distributed proxy servers, called \textbf{Edge Nodes} or \textbf{Points of Presence} (PoPs), deployed inside local ISP networks.

CDNs rely on three core pillars:

\begin{itemize}
\item \textbf{Edge caching:} stateless HTTP video chunks are stored close to end users.
\item \textbf{Request routing:} intelligent DNS or Anycast redirects clients to the closest edge node.
\item \textbf{Server consolidation:} origin server load is reduced while global traffic spikes are absorbed by edge nodes.
\end{itemize}

\begin{center}\includegraphics[width=0.968\linewidth]{Pasted_image_20260520134319.png}\end{center}
\textit{Figure 14: CDN architecture with an origin server synchronizing to European and US edge nodes; local users receive content from nearby PoPs instead of using a high-latency path to the origin.}

\section{Adaptive Bitrate Streaming}

\subsection{The Architecture of Adaptive Streaming}

The industrial standard is \textbf{MPEG-DASH}, \textit{Dynamic Adaptive Streaming over HTTP}, standardized as ISO/IEC 23009-1. It is the first international standard for HTTP-based adaptive bitrate streaming.

The core architectural shift is:

\begin{itemize}
\item \textbf{Dynamic:} video quality adapts during playback according to real-time network conditions.
\item \textbf{Adaptive:} the adjustment logic is client-driven; intelligence moves to the network edge.
\item \textbf{Streaming over HTTP:} standard web servers and CDNs act as passive, stateless file dispensers.
\end{itemize}

\subsection{Video Segments and Representations Levels}

A video sequence is divided into $N$ segments, each with duration $T_S$.

Segment $n$ is encoded at different resolution/quality combinations, producing $K$ \textbf{representations} or \textbf{levels}. Level $k \in \{1, 2, \ldots, K\}$ is characterized by:

\begin{itemize}
\item coding rate $R_C(k)$;
\item resolution $Res_k$;
\item frame rate $F_k$.
\end{itemize}

These values do not vary from one segment to the next inside the same level, but they vary across levels. The values are listed in the \textbf{manifest file}.

\begin{center}
\begin{tabular}{|l|l|l|l|}
\hline
\textbf{Level} & \textbf{Bitrate} & \textbf{Resolution} & \textbf{Frame rate} \\ \hline
1 & 200 kbps & $270 \times 480$ & 25 fps \\ \hline
2 & 500 kbps & $540 \times 960$ & 50 fps \\ \hline
3 & 1 Mbps & $540 \times 960$ & 50 fps \\ \hline
... & ... & ... & ... \\ \hline
\end{tabular}
\end{center}
\begin{center}\includegraphics[width=0.887\linewidth]{Pasted_image_20260520134348.png}\end{center}
\textit{Figure 15: Grid of video thumbnails showing segments 1, 2, and 3 encoded at level 1 low quality, level 2 medium quality, and level 3 high quality.}

\subsection{The Media Presentation Description MPD Manifest File}

The \textbf{Media Presentation Description} (MPD) is a structured XML document that profiles the multimedia asset hosted on the server.

Its hierarchy is:

\begin{itemize}
\item \textbf{Period:} high-level temporal portions of the timeline, such as main content and advertisement insertion.
\item \textbf{AdaptationSet:} alternative media tracks of the same type, such as video track, English audio, and Italian audio.
\item \textbf{Representation:} one quality tier inside an AdaptationSet, such as 1080p at 5 Mbps or 720p at 2.5 Mbps. It specifies codec, frame rate, resolution, and bitrate.
\item \textbf{Segment:} the lowest logical node, containing exact HTTP URLs or byte-range templates for media chunks.
\end{itemize}

\begin{center}\includegraphics[width=1.000\linewidth]{Pasted_image_20260520134407.png}\end{center}
\textit{Figure 16: MPD hierarchy with periods, video/audio AdaptationSets, representations with $R_C$, frame rate, resolution, and codec, and segment URLs pointing to server-side segment tiles.}

\subsection{Video Segment Generation}

Bitrate adaptation can happen only at each new segment. Therefore:

\begin{itemize}
\item every segment must start with a \textbf{closed GOP}, meaning a GOP beginning with an \textbf{IDR frame};
\item open GOPs beginning with \textbf{CRA} frames are forbidden at segment boundaries;
\item no frame inside Segment $N+1$ can reference pixels from Segment $N$.
\end{itemize}

\textbf{Cross-representation alignment} is mandatory: segments across all quality representations, such as 360p, 720p, and 1080p, must have the same temporal duration $T_S$, for example exactly $2.00$ seconds, and identical IDR timestamps.

Segment duration creates a trade-off:

\begin{center}
\begin{tabular}{|l|l|l|}
\hline
\textbf{Segment duration} & \textbf{Advantages} & \textbf{Disadvantages} \\ \hline
Short, 1-2 s & High granularity and fast adaptation to network drops & Lower compression efficiency because IDR frames are frequent and heavy \\ \hline
Long, 6-10 s & Better compression efficiency with fewer IDR overheads & Slower reaction to network changes and higher starvation risk \\ \hline
\end{tabular}
\end{center}
\begin{center}\includegraphics[width=0.968\linewidth]{Pasted_image_20260520134429.png}\end{center}
\textit{Figure 17: Aligned timelines for 1080p, 720p, and 360p tiers, each split into 2 s segments with closed GOPs starting with IDR frames at boundaries and open GOPs with CRA only inside segments.}

\subsection{Live Streaming vs Stored Video Paradigms}

\textbf{Stored video}, or video on demand, is asynchronous with transmission. Media is pre-encoded, pre-segmented, and cached on CDNs in advance; latency is not critical.

\textbf{Live streaming} is simultaneous and contextual to transmission. Segments are created dynamically by a real-time encoder.

DASH fits live streaming because live audiences tolerate relaxed delays compared with conversational WebRTC. A 10-30 s broadcast lag behind the physical event is acceptable, so DASH can preserve scalability and caching advantages.

\subsection{Low-Latency Streaming Frameworks}

Traditional DASH/HLS latency is tied to segment length: a server cannot publish or cache a 6 s segment until encoding of the full segment is complete.

\textbf{CMAF}, the \textit{Common Media Application Format}, breaks a structural segment into \textbf{CMAF chunks} of about 200-500 ms. Chunks do not require independent IDR frames; they share the parent segment's anchor IDR.

\textbf{HTTP Chunked Transfer Encoding} (CTE) enables pipelined delivery through HTTP/1.1 and HTTP/2. As soon as the encoder finishes a 200 ms chunk, the origin server streams it through the CDN pipeline before the parent segment is closed.

The target is \textbf{1-3 s glass-to-glass latency}, comparable to traditional cable TV/broadcast, without abandoning commodity HTTP caching.

\begin{center}\includegraphics[width=0.968\linewidth]{Pasted_image_20260520134456.png}\end{center}
\textit{Figure 18: Traditional DASH latency diagram where a 6 s segment must be completely encoded before publishing to CDN and player.}

Figure 19: Low-latency DASH/CMAF diagram where a 6 s parent segment contains 200 ms chunks streamed immediately to CDN and player.

\subsection{Network Metrics Throughput and Capacity}

\textbf{Link bitrate} $R$ is the transmission rate supported by the physical-layer interface, measured in bits per second.

\begin{calloutbox}{red!70!black}{Instantaneous Throughput}
\textbf{Statement:} \textbf{Instantaneous throughput} $S(t)$ is the actual time-varying rate at which data bytes successfully arrive at the application layer:

\[
S(t) = \frac{dD(t)}{dt} \leq R
\]

where $D(t)$ is cumulative data successfully received up to time $t$.

\textbf{Interpretation:} throughput is lower than physical link bitrate because of headers, losses, flow control, and congestion control.
\end{calloutbox}

\begin{calloutbox}{red!70!black}{Average Throughput}
\textbf{Statement:} \textbf{Average throughput} over a window $T$ smooths network bursts and losses:

\[
\bar{S} = \frac{1}{T} \int_0^T S(t)\,dt = \frac{D(T)}{T}
\]

\textbf{Interpretation:} ABR algorithms often use smoothed throughput estimates rather than raw instantaneous samples.
\end{calloutbox}

\subsection{Nodal Delay and End-to-End Delay}

The total end-to-end nodal delay $d_{e2e}$ of a packet is the sum of delay factors at each network hop $k$:

\begin{calloutbox}{red!70!black}{End-to-End Nodal Delay}
\textbf{Statement:}

\[
d_{e2e} = \sum_k d_{\text{proc}}(k) + d_{\text{queue}}(k) + d_{\text{tx}}(k) + d_{\text{prop}}(k)
\]

\textbf{Interpretation:} total delay depends on processing, queuing, transmission, and propagation components at every hop.
\end{calloutbox}

The components are:

\begin{itemize}
\item \textbf{Processing delay} $d_{\text{proc}}(k)$;
\item \textbf{Queuing delay} $d_{\text{queue}}(k)$;
\item \textbf{Transmission delay} $d_{\text{tx}}(k)$;
\item \textbf{Propagation delay} $d_{\text{prop}}(k)$.
\end{itemize}

\subsubsection{Nodal Delay Breakdown}

\begin{itemize}
\item \textbf{Processing delay} $d_{\text{proc}}(k)$: time routers spend checking bit-level errors, reading packet headers, and deciding the output link.
\item \textbf{Queuing delay} $d_{\text{queue}}(k)$: time spent waiting in router buffers due to congestion. It is dynamic and highly variable.
\item \textbf{Transmission delay} $d_{\text{tx}}(k)$: time needed to push all packet bits into the physical link:
\end{itemize}

\[
d_{\text{tx}}(k) = \frac{L}{R}
\]

where $L$ is packet length and $R$ is link rate.

\begin{itemize}
\item \textbf{Propagation delay} $d_{\text{prop}}(k)$: travel time through the physical medium:
\end{itemize}

\[
d_{\text{prop}}(k) = \frac{d}{s}
\]

where $d$ is distance and $s$ is propagation speed.

Increasing $R$ reduces $d_{\text{tx}}$, but does not reduce $d_{\text{prop}}$ or $d_{\text{queue}}$. Distance and congestion remain structural bottlenecks for live interactive video.

\subsection{Jitter}

Packet delay is caused by many different phenomena. We care about average delay and its \textbf{variance}, called \textbf{jitter}.

High jitter means that delay can be much different from its average. The system must cope with occasional peaks in delay.

\begin{center}\includegraphics[width=0.806\linewidth]{Pasted_image_20260520134551.png}\end{center}
\textit{Figure 20: Histogram of per-packet delay in an end-to-end Internet connection, with average, minimum, and maximum delay markers; most packets are below 120 ms, while some reach about 150 ms.}

\subsection{The Playout Buffer The Client Defensive Reservoir}

To absorb delay variability, the client uses a \textbf{playout buffer}. It isolates stochastic packet arrivals from deterministic frame decoding.

The main state variable is buffer size $B(t)$, measured in seconds. It represents how much video has been received but not played yet. If the network fails at time $t$, playout can continue for another $B(t)$ seconds.

\subsection{Playout Buffer Dynamics}

The model is continuous-time and continuous-valued, using a \textbf{fluid dynamics framework}.

At the beginning, the buffer is empty. When the first bits of the first segment are received, playout usually does not start immediately. The client first builds a reservoir of $L$ segments. When:

\[
B(t) = L \cdot T_S
\]

playout starts. The buffer can then grow or shrink depending on network conditions and requested coding rate $R_C$.

\textbf{Buffer starvation} occurs when:

\[
B(t) = 0
\]

The tank is empty and playout freezes, producing a rebuffering stall. Playout resumes only when $M$ segments have been received, so after a stall:

\[
B(t) = M \cdot T_S
\]

Typically $L \geq M$, because users tolerate initial delay more than mid-video stalls.

\subsubsection{Fluid Model}

The buffer dynamics are modeled using input and output flow rates:

\[
\frac{dB}{dt} = f_{\text{in}} - f_{\text{out}}
\]

\textbf{Input flow rate} $f_{\text{in}}$ is the number of seconds of video received per real-time second. If the client receives $\bar{S}T$ bits over an interval of duration $T$, and the coding rate is $R_C$, the amount of playable video received is:

\[
D(T) = \frac{\bar{S}T}{R_C}
\]

Thus the average input flow rate is:

\[
\frac{D(T)}{T} = \frac{\bar{S}}{R_C}
\]

and the instantaneous input flow rate is:

\[
\lim_{T \to 0} \frac{D(T)}{T} = \frac{S(t)}{R_C}
\]

This input is driven by network conditions.

\textbf{Output flow rate} $f_{\text{out}}$ is the number of seconds of video consumed per real-time second. During active playback $f_{\text{out}} = 1$. During rebuffering $f_{\text{out}} = 0$.

\begin{calloutbox}{red!70!black}{Playout Buffer Differential Equation}
\textbf{Statement:}

\[
\frac{dB}{dt} =
\begin{cases}
\dfrac{S}{R_C} - 1 & \text{if playout} \\
\dfrac{S}{R_C} & \text{if rebuffering}
\end{cases}
\]

\textbf{Interpretation:} buffer grows when data arrives faster than video is consumed, shrinks when coding rate exceeds available throughput during playback, and only fills during rebuffering.
\end{calloutbox}

\subsection{Examples}

\begin{calloutbox}{green!55!black}{Sudden Throughput Drop in ABR}
\textbf{Problem:} Consider an ABR system with $L = 2$, $M = 1$, and $N = 6$, where throughput suddenly drops during segment download.

\textbf{Approach:} Track segment download and playout timelines and observe $B(t)$.

\textbf{Solution:} Segments 1-3 download quickly. Segment 4 takes longer because throughput drops. During this period the buffer decreases, reaches zero, and causes rebuffering. After recovery, playout continues with segments 5 and 6.

\textbf{Result:} the buffer rises to about $2T_S$, declines through $T_S$, reaches $0$, triggers rebuffering, then recovers.

\textbf{Takeaway:} sudden throughput drops can cause starvation even if previous segments were downloaded quickly.
\end{calloutbox}

\begin{center}\includegraphics[width=1.000\linewidth]{Pasted_image_20260520134648.png}\end{center}
\textit{Figure 21: Timeline showing download and playout of segments 1-6, a throughput drop during segment 4, a red buffer curve reaching zero, and rebuffering before playback continues.}

\subsubsection{Interactive Matlab Simulations}

All simulations use a network throughput:

\begin{itemize}
\item $S_1 = 3$ Mbps for $0$-$10$ s;
\item $S_2 = 1$ Mbps for $10$-$40$ s;
\item $S_3 = 3$ Mbps after $40$ s;
\item $R_C = 2.5$ Mbps.
\end{itemize}

\begin{calloutbox}{green!55!black}{Matlab Simulation 1}
\textbf{Problem:} $T_S = 0.5$, $L = 2$, $M = 1$.

\textbf{Solution:} Small initial buffer capacity gives fast setup but frequent short stalls.

\textbf{Result:} number of stalls: 12. Total stall duration: 15.00 s.

\textbf{Takeaway:} low $L$ and low $M$ reduce initial delay but produce many annoying stalls.
\end{calloutbox}

\begin{center}\includegraphics[width=0.968\linewidth]{Pasted_image_20260520134714.png}\end{center}
\textit{Figure 22: Matlab plots showing throughput $S(t)$ versus coding rate $R_C$ and buffer $B(t)$ oscillating near zero with many rebuffering events.}

\begin{calloutbox}{green!55!black}{Matlab Simulation 2}
\textbf{Problem:} $T_S = 0.5$, $L = 10$, $M = 1$.

\textbf{Solution:} Larger initial buffer capacity reduces stalls but keeps roughly the same stall frequency pattern once the buffer falls low.

\textbf{Result:} number of stalls: 10. Total stall duration: 11.95 s.

\textbf{Takeaway:} increasing only $L$ helps, but if $M$ remains low, recovery after each stall is weak.
\end{calloutbox}

\begin{center}\includegraphics[width=0.968\linewidth]{Pasted_image_20260520134728.png}\end{center}
\textit{Figure 23: Matlab plots showing larger initial buffer rise followed by oscillations similar to the first case.}

\begin{calloutbox}{green!55!black}{Matlab Simulation 3}
\textbf{Problem:} $T_S = 0.5$, $L = 10$, $M = 5$.

\textbf{Solution:} Larger initial buffer and larger post-stall refill reduce the number of stalls but make stalls longer.

\textbf{Result:} number of stalls: 2. Total stall duration: 12.50 s.

\textbf{Takeaway:} increasing $M$ trades many short stalls for fewer longer stalls.
\end{calloutbox}

\begin{center}\includegraphics[width=0.968\linewidth]{Pasted_image_20260520134744.png}\end{center}
\textit{Figure 24: Matlab plots showing fewer but longer rebuffering events.}

\begin{calloutbox}{green!55!black}{Matlab Simulation 4}
\textbf{Problem:} $T_S = 2$, $L = 10$, $M = 5$.

\textbf{Solution:} Very large initial buffer capacity absorbs the capacity drop completely.

\textbf{Result:} number of stalls: 0. Total stall duration: 0.00 s.

\textbf{Takeaway:} larger $T_S$ and buffer reserve can eliminate stalls, but increase latency and reduce adaptation granularity.
\end{calloutbox}

\begin{center}\includegraphics[width=0.968\linewidth]{Pasted_image_20260520134757.png}\end{center}
\textit{Figure 25: Matlab plots showing buffer rising to about 20 s and no rebuffering events.}

\begin{calloutbox}{green!55!black}{Matlab Simulation 5}
\textbf{Problem:} $T_S = 0.5$, $L = 2$, $M = 1$, repeated for reference.

\textbf{Solution:} This repeats the first simulation configuration.

\textbf{Result:} same behavior as Matlab Simulation 1.

\textbf{Takeaway:} reference case for comparing the effect of changing $L$, $M$, and $T_S$.
\end{calloutbox}

\begin{center}\includegraphics[width=0.968\linewidth]{Pasted_image_20260520134810.png}\end{center}
\textit{Figure 26: Same Matlab simulation plot as the first configuration.}

\subsection{Wrap-Up System Model for AS}

The main parameters of the adaptive streaming model are:

\begin{center}
\begin{tabular}{|l|l|}
\hline
\textbf{Parameter} & \textbf{Meaning} \\ \hline
$T_S$ & Playout duration of a video segment, assumed constant \\ \hline
$S_n$ & Throughput during download of segment $n$ \\ \hline
$q(n)$ & Requested level for segment $n$ \\ \hline
$T_D(n,k)$ & Downloading time for segment $n$ at level $k$ \\ \hline
$B(t)$ & Playout buffer, in seconds, meaning received but not yet played video \\ \hline
$T_{IB}$ & Initial buffering time \\ \hline
$M \cdot T_S$ & Rebuffering constraint during playout \\ \hline
\end{tabular}
\end{center}

The download time is:

\[
T_D(n,k) = \frac{T_S R_C(k)}{S_n}
\]

The initial buffering time is:

\[
T_{IB} = L \cdot T_S
\]

before playback starts, the client fills the playout buffer with $L$ segments.

\section{Video Streaming QoE}

\subsection{Video Streaming Quality of Experience}

The problem is to define \textbf{Quality of Experience} (QoE) for a streaming service.

Global QoE is influenced by:

\begin{enumerate}
\item per-segment video quality;
\item number and duration of rebuffering events;
\item number and amplitude of quality changes;
\item duration of initial buffering time.
\end{enumerate}

\subsubsection{Per-Segment Video Quality}

For segment $n$, video quality is generally:

\[
Q(n) = f(R(k_n))
\]

where $n$ is the segment index and $k_n$ is the selected level.

The function $f$ should be a rate-quality function, such as:

\[
\text{MOS} = f(R)
\]

However, perceived quality versus rate depends strongly on content. Some images may look perfect at 0.5 bpp, while others show visible artifacts at 1 bpp. PSNR is only a proxy for subjective quality.

In relevant cases, perceived quality is a monotonic increasing function of coding rate $R$. Since modeling MOS versus coding rate is difficult, simple models are often used:

\[
Q(n) = R(k_n)
\]

or even:

\[
Q(n) = k_n
\]

\subsubsection{Rebuffering Events}

Let $\Delta_n$ be the duration of the rebuffering event that occurs during playout of segment $n$. If no rebuffering affects segment $n$, then:

\[
\Delta_n = 0
\]

After rebuffering, at least one full segment is downloaded before resuming playout, so there cannot be more than one rebuffering event during a single segment playout.

Rebuffering is highly disruptive for QoE, regardless of duration. QoE during segment $n$ decreases with stall duration $\Delta_n$. This reduction is modeled by a function $\phi(\Delta)$ such that $\phi(0)=0$, while $\phi(0^+)$ is relatively large and then increases slowly:

\begin{calloutbox}{red!70!black}{Rebuffering Penalty}
\textbf{Statement:}

\[
\phi(\Delta) =
\begin{cases}
0 & \text{if } \Delta = 0 \\
a\Delta + b & \text{if } \Delta > 0
\end{cases}
\]

where $a$ is small and $b$ is large.

\textbf{Interpretation:} even a tiny stall creates a large QoE penalty, while longer stalls add penalty more slowly.
\end{calloutbox}

\begin{center}\includegraphics[width=0.806\linewidth]{Pasted_image_20260520134841.png}\end{center}
\textit{Figure 27: Plot of $\phi$ versus $\Delta_n$ with a discontinuous jump at zero to $b$, then a linear increase with slope $a$.}

\subsubsection{Quality Changes and Initial Buffering}

Any quality variation, both decrease and increase, is annoying by itself. A constant but slightly lower quality can be preferable to frequent level changes.

Users usually tolerate some initial buffering if it reduces later rebuffering. However, excessive initial buffering also reduces QoE.

\subsection{A Video Quality Metric for Streaming}

\begin{calloutbox}{red!70!black}{Segment-Level QoE Metric}
\textbf{Statement:} QoE for segment $n$ can be represented as:

\[
J(n) = \lambda_1 k_n - \lambda_2 |k_n - k_{n-1}| - \phi(\Delta_n)
\]

\textbf{Interpretation:} $J(n)$ increases with selected level $k_n$, decreases when quality changes from the previous segment, and decreases sharply when a stall occurs.
\end{calloutbox}

The sequence-level QoE is:

\begin{calloutbox}{red!70!black}{Sequence-Level QoE}
\textbf{Statement:}

\[
J = \sum_{n=1}^{N} J(n) - \lambda_3 T_{ST}
\]

where $T_{ST}$ is the time required to download $L$ segments with the quality level assigned by the streaming strategy.

\textbf{Interpretation:} total QoE sums segment utility and subtracts annoyance from initial buffering.
\end{calloutbox}

It is difficult to set the values of $\lambda_i$ and the behavior of $\phi$. Usually $\lambda_3$ is relatively small, because users tolerate initial delay, while $\phi$ has a discontinuity at $0$ because any buffering event is annoying.

\section{Adaptive Bitrate ABR Strategies}

\subsection{Adaptive Bit Rate ABR}

A client uses an \textbf{Adaptive Bit Rate} (ABR) algorithm to select the best quality level for the next segment that can be downloaded in time for playback.

ABR is typically not standardized because it is implemented only on the client side. Common ABR algorithms are:

\begin{itemize}
\item \textbf{throughput-based};
\item \textbf{buffer-based};
\item \textbf{hybrid}.
\end{itemize}

\subsection{Throughput-Based ABR}

Throughput-based schemes estimate network throughput and decide the bitrate for the next segment:

\begin{itemize}
\item choose quality so the estimated download time is not longer than segment playout time;
\item examples include PANDA, FESTIVE, and SQUAD.
\end{itemize}

References:

\begin{itemize}
\item Z. Li, X. Zhu, J. Gahm, R. Pan, H. Hu, A. C. Begen, and D. Oran, "Probe and adapt: Rate adaptation for HTTP video streaming at scale," \textit{IEEE Journal on Selected Areas in Communications}, vol. 32, no. 4, 2014.
\item J. Jiang, V. Sekar, and H. Zhang, "Improving fairness, efficiency, and stability in HTTP-based adaptive video streaming with FESTIVE," in \textit{Proceedings of the Conference on Emerging Networking Experiments and Technologies}, Association for Computing Machinery, 2012.
\end{itemize}

\subsection{Buffer-Based ABR}

Buffer-based algorithms use buffer level to decide the bitrate for the next segment:

\begin{itemize}
\item if buffer level is high, frames are arriving faster than consumed, so quality can be increased;
\item if buffer level is low, frames arrive slower than consumed, so quality should be decreased.
\end{itemize}

Examples include BBA and \textbf{BOLA}.

References:

\begin{itemize}
\item T.-Y. Huang, R. Johari, N. McKeown, M. Trunnell, and M. Watson, "A buffer-based approach to rate adaptation: Evidence from a large video streaming service," in \textit{Proceedings of the 2014 ACM Conference on SIGCOMM}, Association for Computing Machinery, 2014.
\item K. Spiteri, R. Urgaonkar, and R. K. Sitaraman, "BOLA: Near-optimal bitrate adaptation for online videos," in \textit{INFOCOM}, IEEE, 2016.
\end{itemize}

\subsection{Hybrid ABR}

Hybrid algorithms use both throughput prediction and buffer level. Examples include ELASTIC, MPC, ABMA+, and DYNAMIC.

References:

\begin{itemize}
\item X. Yin, A. Jindal, V. Sekar, and B. Sinopoli, "A control-theoretic approach for dynamic adaptive video streaming over HTTP," \textit{SIGCOMM Computer Communication Review}, vol. 45, no. 4, 2015.
\item A. Beben et al., "ABMA+ lightweight and efficient algorithm for HTTP adaptive streaming," \textit{Proceedings of the 7th International Conference on Multimedia Systems}, 2016.
\end{itemize}

\subsection{An Example of ABR Adaptive Bitrate Streaming}

The client chooses $k_n$ with the goal of maximizing:

\[
J = \lambda_1 k_n - \lambda_2 |k_n - k_{n-1}| - \phi(\Delta_n)
\]

The client must decide the quality level for segment $n$, knowing:

\begin{itemize}
\item previous level $k_{n-1}$;
\item playout buffer size $B_0 > 0$, meaning $B_0$ seconds of video have been received but not played;
\item throughput estimate $\hat{S}$;
\item segment duration $T_S$.
\end{itemize}

The algorithm in the source is:

\begin{verbatim}
Set J_min = infinity

For each ell = 1 ... K:
    Set R = R(ell), the coding rate at level ell
    Set beta = S_hat / R - 1
    Assuming constant throughput during the download of the segment,
    B(t) = B_0 + beta t

    If beta >= 0:
        Delta = 0
        evaluate J

    If beta < 0:
        the buffer starts depleting during the download of segment n

\end{verbatim}

If $\beta < 0$, the buffer reaches zero at time $t_0$ such that:

\[
B_0 + \beta t_0 = 0
\]

therefore:

\[
t_0 =
\frac{B_0}{1 - \frac{\hat{S}}{R}}
=
\frac{R B_0}{R - \hat{S}}
\]

If $t_0 > T_S$, playout of the new segment finishes before the buffer is empty, so:

\[
\Delta = 0
\]

Otherwise, if $t_0 < T_S$, the buffer reaches zero while there are:

\[
T_R = T_S - \frac{R B_0}{R - \hat{S}}
\]

seconds of the segment still to download. A stall occurs. The duration of the stall is the time needed to download the rest of the segment plus $M$ new segments:

\[
\Delta = \frac{(T_R + M T_S)R}{\hat{S}}
\]

Then evaluate:

\[
J = \lambda_1 \ell - \lambda_2 |\ell - k_{n-1}| - \phi(\Delta)
\]

If $J < J_{\min}$, set $k_n = \ell$ and $J_{\min} = J$.

At the end of the loop, $k_n$ is the best quality level under these hypotheses.

This is only one possible ABR strategy. Its simplifying assumptions can undermine behavior:

\begin{itemize}
\item it assumes constant throughput during the full segment download;
\item it assumes constant throughput during rebuffering if a stall happens;
\item it assumes constant coding rate during playout, although intra frames have much higher rates than other frames;
\item it has no safety threshold for buffer level $B$: as long as $B$ does not reach zero during the current segment, it does not act to prevent future stalls.
\end{itemize}

The algorithm runs only on the client side. Each client can run a different strategy, such as aggressive high-quality choices with higher stall risk, or cautious choices with lower stall probability and lower average quality.

\subsection{Limits of Current AS Technology}

Current adaptive streaming technology has several limits:

\begin{itemize}
\item large RTT increases freezing/low-quality risk;
\item larger RTT makes throughput estimation harder;
\item decisions are entirely client-side, so the client has only a partial view of channel conditions;
\item adaptation can be suboptimal with multiple DASH clients or aggressive cross-traffic;
\item size/bitrate is not the only factor affecting perceived quality.
\end{itemize}

\section{Exercises}

\subsection{Exercise 1}

\begin{calloutbox}{green!55!black}{Exercise 1}
\textbf{Problem:} Consider an adaptive streaming system with:

\begin{itemize}
\item segment duration $T_S = 1$ s;
\item number of quality levels $K = 3$;
\item bitrates $R_C(1) = 500\,\text{kbps}$, $R_C(2) = 1\,\text{Mbps}$, $R_C(3) = 2\,\text{Mbps}$.
\end{itemize}

The initial buffering stops after $t_{\text{ST}}$ seconds. This is the duration of initial buffering, not the amount of downloaded playable video.

Set $t=0$ when the first bit is received. The server-client connection has throughput:

\[
\forall t \in (0,T),\quad S(t) = S_1 = 0.4\,\text{Mbps}
\]

\[
\forall t > T,\quad S(t) = S_2 = 0.5\,\text{Mbps}
\]

Find the minimum initial buffering that assures no stalls when $q(n)=1$, meaning the client always requests quality level 1.

\textbf{Approach:} Compare buffer slopes during rebuffering and playout before and after $T$.

\textbf{Solution:} For $t<T$, $S < R_C$. For $t>T$, $S = R_C$. During rebuffering:

\[
B'_{\text{rebuff}}(t) = \frac{S}{R_C}
\]

During playout:

\[
B'_{\text{playout}}(t) = \frac{S}{R_C} - 1
\]

If $t_{\text{ST}} > T$, the slope of the buffer is:

\[
B'(t) =
\begin{cases}
\dfrac{S_1}{R_C} = 0.8 & \forall t \in (0,T) \\
\dfrac{S_2}{R_C} = 1 & \forall t \in (T,t_{\text{ST}}) \\
\dfrac{S_2}{R_C} - 1 = 0 & \forall t > t_{\text{ST}}
\end{cases}
\]

In this case, the buffer never depletes because the slope is positive or zero. There is no rebuffering.

If $t_{\text{ST}} = T$, the buffer reaches $0.8T$ before playout starts, then stays constant.

If $t_{\text{ST}} < T$, during $(0,t_{\text{ST}})$ the buffer fills with slope:

\[
B' = \frac{S_1}{R_C} = 0.8
\]

so:

\[
B(t_{\text{ST}}) = 0.8\,t_{\text{ST}}
\]

For $t > t_{\text{ST}}$ until the buffer is empty, the slope is:

\[
B' = \frac{S_1}{R_C} - 1 = -0.2
\]

hence:

\[
B(t) = B(t_{\text{ST}}) + B'(t-t_{\text{ST}})
= 0.8\,t_{\text{ST}} - 0.2(t-t_{\text{ST}})
= t_{\text{ST}} - 0.2t
\]

The buffer decreases linearly until it reaches zero or until $t=T$, when the slope becomes zero. The zero time $t^*$ is:

\[
B(t^*) = 0
\Leftrightarrow t_{\text{ST}} - 0.2t^* = 0
\Leftrightarrow t^* = 5t_{\text{ST}}
\]

To avoid rebuffering, the buffer must not empty before the slope changes at $T$:

\[
T < 5t_{\text{ST}}
\]

The minimum is obtained at:

\[
T = 5t_{\text{ST}}
\]

therefore:

\[
t_{\text{ST}} = \frac{T}{5}
\]

\textbf{Result:} the minimum initial buffering is $t_{\text{ST}} = T/5$.

\textbf{Takeaway:} initial buffering must be long enough for the buffer to survive the low-throughput interval before capacity becomes equal to the selected coding rate.
\end{calloutbox}

\begin{center}\includegraphics[width=0.565\linewidth]{Pasted_image_20260520134935.png}\end{center}
\textit{Figure 28: Two plots of $B(t)$ for $t_{\text{ST}}>T$ and $t_{\text{ST}}=T$, showing buffer growth before playback and no depletion.}

\begin{center}\includegraphics[width=1.000\linewidth]{Pasted_image_20260520134955.png}\end{center}
\textit{Figure 29: Two plots for $t_{\text{ST}}<T$, with one case where the buffer reaches zero at $t^*=5t_{\text{ST}}$ and the minimum case $t_{\text{ST}}=T/5$ where the buffer reaches zero exactly at $T$.}

\subsection{Exercise 2}

An adaptive streaming system has:

\begin{itemize}
\item segment duration $T_S = 1$ s;
\item number of quality levels $K = 3$;
\item bitrates $R_C(1)=500\,\text{kbps}$, $R_C(2)=1\,\text{Mbps}$, $R_C(3)=2\,\text{Mbps}$;
\item initial buffering stops after downloading $L=2$ segments;
\item after rebuffering, playout resumes when the next segment is loaded, so $M=1$.
\end{itemize}

Set $t=0$ when the first bit is received. Throughput is:

\[
S(t) =
\begin{cases}
S_1 = 1.5\,\text{Mbps} & \forall t \in (0,2) \\
S_2 = 0.1\,\text{Mbps} & \forall t \in (2,3) \\
S_3 = 1.8\,\text{Mbps} & \forall t > 3
\end{cases}
\]

Questions:

\begin{enumerate}
\item If the client always requests maximum quality, $q(n)=3$ for all $n \geq 1$, when does playout begin?
\item With $q(n)=3$ for all $n \geq 1$, what are the frequency and duration of rebuffering, if any?
\item Show that selecting quality according to
\end{enumerate}

\[
q: n \geq 1 \to
\begin{cases}
3 & \text{if } n \bmod 5 \neq 3 \\
2 & \text{if } n \bmod 5 = 3
\end{cases}
\]

avoids all rebuffering.

\begin{enumerate}
\item Is the previous strategy better than using $q=2$ all the time?
\end{enumerate}

\begin{center}\includegraphics[width=0.806\linewidth]{Pasted_image_20260520135034.png}\end{center}
\textit{Figure 30: Throughput plot $S(t)$ in Mbps versus time, with $S=1.5$ for $t \in (0,2)$, $S=0.1$ for $t \in (2,3)$, and $S=1.8$ for $t>3$.}

\subsubsection{Solution Exercise 2 Question 1}

\begin{calloutbox}{green!55!black}{Exercise 2 Question 1}
\textbf{Problem:} Determine when playout starts if $q(n)=3$ for all $n$.

\textbf{Approach:} Initial buffering ends after $L=2$ segments have been downloaded.

\textbf{Solution:} The number of bits needed is:

\[
N_b = T_S R_C(q(1)) + T_S R_C(q(2))
\]

More generally:

\[
N_b = T_S \sum_{n=1}^{L} R_C(q(n))
\]

Here $q(1)=q(2)=3$, and $R_C(3)=2\,\text{Mbps}$. Each segment takes 2 Mbits, so:

\[
N_b = 2 + 2 = 4\,\text{Mbits}
\]

In the first 2 s:

\[
N_1 = 2 \cdot 1.5 = 3\,\text{Mbits}
\]

In the next 1 s:

\[
N_2 = 1 \cdot 0.1 = 0.1\,\text{Mbits}
\]

Remaining bits:

\[
N_b - N_1 - N_2 = 0.9\,\text{Mbits}
\]

Starting from $t=3$, throughput is $S_3=1.8\,\text{Mbps}$, so the remaining download time is:

\[
\frac{0.9}{1.8} = 0.5\,\text{s}
\]

\textbf{Result:} playout starts at:

\[
t_{\text{ST}} = 3.5\,\text{s}
\]

\textbf{Takeaway:} startup time depends on the cumulative bits required for the first $L$ requested segments and the piecewise throughput profile.
\end{calloutbox}

\subsubsection{Solution Exercise 2 Question 2}

\begin{calloutbox}{green!55!black}{Exercise 2 Question 2}
\textbf{Problem:} Find frequency and duration of rebuffering for $q(n)=3$.

\textbf{Approach:} After $t_{\text{ST}}=3.5$, throughput is constant at $S_3=1.8\,\text{Mbps}$ and $R_C=2\,\text{Mbps}$.

\textbf{Solution:} Since $R_C>S$, the buffer decreases. During playout:

\[
B'(t) = \frac{S_3}{R_C} - 1 = \frac{1.8}{2} - 1 = -\frac{1}{10}
\]

Starting from $B(t_{\text{ST}})=2$, the buffer empties after:

\[
\frac{B(t_{\text{ST}})}{|B'|} = \frac{2}{1/10} = 20\,\text{s}
\]

During these 20 s of playout, the client receives:

\[
20 \cdot 1.8 = 36\,\text{Mbits}
\]

corresponding to:

\[
\frac{36}{R_C} = \frac{36}{2} = 18\,\text{s}
\]

of playable video. Together with the initial 2 s, received and played video equals $2+18=20$ s, so the buffer becomes empty at:

\[
t_{\text{rebuff}} = 3.5 + 20 = 23.5\,\text{s}
\]

Rebuffering then starts. The current segment has been fully downloaded, so only the next segment is needed:

\[
N_{\text{rebuff}} = T_S R_C(q(n)) = 2\,\text{Mbits}
\]

The source states this takes:

\[
\frac{N_{\text{rebuff}}}{S} = \frac{2}{1.5} = \frac{10}{9} \approx 1.1\overline{1}\,\text{s}
\]

(note: possible OCR artifact, because for $t>3$ the stated throughput profile gives $S_3=1.8\,\text{Mbps}$, which would produce $2/1.8=10/9$ s; the final numerical value $10/9$ matches $S=1.8$ rather than $1.5$. The original expression is preserved.)

After rebuffering, the buffer has 1 s of playable video. It empties again after 10 s, then the next rebuffering lasts again $10/9$ s.

\textbf{Result:} the rebuffering period is:

\[
T_{\text{rebuff}} = 11.\overline{1}\,\text{s}
\]

and stall duration is:

\[
\Delta = \frac{10}{9}\,\text{s}
\]

\textbf{Takeaway:} because $S_3<R_C(3)$, maximum quality creates periodic rebuffering after startup.
\end{calloutbox}

\begin{center}\includegraphics[width=0.726\linewidth]{Pasted_image_20260520135123.png}\end{center}
\textit{Figure 31: Buffer plot showing initial rise to $B=2$ at $t_{\text{ST}}=3.5$, slow decrease to zero at about 23.5 s, and periodic rebuffering every about 11 s with buffer oscillating between 0 and about 1 s.}

\subsubsection{Solution Exercise 2 Question 3}

\begin{calloutbox}{green!55!black}{Exercise 2 Question 3}
\textbf{Problem:} Show that the periodic strategy using level 2 every fifth segment avoids rebuffering.

\textbf{Approach:} Analyze one 5-segment cycle after startup.

\textbf{Solution:} Initial rebuffering ends at $t=3.5$ s because the first two segments are requested with $k=3$, as before. Therefore:

\[
B(3.5)=2
\]

The third segment is requested at $k=2$. Its download time is:

\[
T_D(3) = \frac{T_S R_C(2)}{S_3} = \frac{1}{1.8}\,\text{s} = \frac{5}{9}\,\text{s}
\]

During this time, the buffer slope is:

\[
B'(t)=\frac{S}{R_C}-1 = 0.8\,\text{s/s}
\]

The slope is positive, so the buffer fills. After downloading the segment:

\[
B = B(3.5) + B'(t)T_D(3)
= 2 + \frac{5}{9}\cdot\frac{4}{5}
= 2.\overline{4}\,\text{s}
\]

Now four segments, from 4 to 7, are downloaded at level $k=3$. Since $S_n$ is constant, download time is:

\[
T_D = 4T_D(n,3)
= 4 \frac{T_S R_C(3)}{S_n}
= 4 \cdot \frac{2}{1.8}\,\text{s}
= \frac{40}{9}\,\text{s}
= 4.\overline{4}\,\text{s}
\]

During these four downloads the buffer is depleted. When segment 7 finishes, the buffer is:

\[
B(t) = 2.\overline{4} - \frac{4.\overline{4}}{10} = 2\,\text{s}
\]

The system has returned to the same buffer level at the start of the cycle, so the pattern repeats every 5 segments.

\textbf{Result:} the buffer oscillates between $2$ and $2.\overline{4}$ seconds and never reaches zero.

\textbf{Takeaway:} periodic quality reduction provides enough buffer gain to compensate for four high-quality segments.
\end{calloutbox}

\begin{center}\includegraphics[width=0.887\linewidth]{Pasted_image_20260520135147.png}\end{center}
\textit{Figure 32: Stem plot of $q(n)$ from $n=1$ to $15$, with $q=3$ for most segments and $q=2$ at $n=3,8,13$.}

\begin{center}\includegraphics[width=1.000\linewidth]{Pasted_image_20260520135203.png}\end{center}
Figure 33: Buffer plot from 4 to 14 s showing periodic rise from 2 to $2.\overline{4}$ during level-2 segment download and slow fall back to 2 during four level-3 segment downloads.

\subsubsection{Exercise 2 Question 4}

Compare two strategies:

\[
q_1: n \geq 1 \to
\begin{cases}
3 & \text{if } n \bmod 5 \neq 3 \\
2 & \text{if } n \bmod 5 = 3
\end{cases}
\]

and:

\[
q_2: n \geq 1 \to 2 \quad \forall n \geq 1
\]

using:

\[
J = \sum_{n=1}^{N} J(n) - \lambda_3 T_{ST}
\]

and:

\[
J(n) = \lambda_1 k_n - \lambda_2 |k_n-k_{n-1}| - \phi(\Delta_n)
\]

\begin{calloutbox}{green!55!black}{Exercise 2 Question 4}
\textbf{Problem:} Decide whether periodic switching is better than always using $q=2$.

\textbf{Approach:} Compare QoE for the two no-stall strategies.

\textbf{Solution:} For the first strategy, there is no stall, so:

\[
J_1 =
\sum_{n=1}^{N}
[\lambda_1 q_1(n) - \lambda_2 |q_1(n)-q_1(n-1)|]
- \lambda_3 T_1
\]

Here $q_1=3$ for $4N/5$ segments and $q_1=2$ for $N/5$ segments. Also, $|q_1(n)-q_1(n-1)|=1$ for $2N/5$ segments. Therefore:

\[
J_1 =
\frac{12}{5}N\lambda_1
+ \frac{2}{5}N\lambda_1
- \frac{2}{5}N\lambda_2
- \lambda_3 T_1
\]

\[
J_1 =
N\left(\frac{14}{5}\lambda_1 - \frac{2}{5}\lambda_2\right)
- \lambda_3 T_1
\]

In the second case there is no stall either. The QoE is:

\[
J_2 =
\sum_{n=1}^{N}
[\lambda_1 q_2(n) - \lambda_2 |q_2(n)-q_2(n-1)|]
- \lambda_3 T_2
\]

Since $q_2=2$ always and there are no changes:

\[
J_2 = 2N\lambda_1 - \lambda_3 T_2
\]

For large $N$:

\[
J_1 - J_2
\approx
N\left(\frac{4}{5}\lambda_1 - \frac{2}{5}\lambda_2\right)
\]

This is positive if and only if:

\[
\lambda_1 > \frac{\lambda_2}{2}
\]

\textbf{Result:} periodic switching is better than constant level 2 only when the reward for higher quality is large enough compared with the penalty for quality changes.

\textbf{Takeaway:} the best strategy depends on subjective weights $\lambda_i$, which are usually evaluated through subjective tests.
\end{calloutbox}

\section{Final Summary Table}

\begin{center}
\begin{tabular}{|l|l|l|l|}
\hline
\textbf{Concept} & \textbf{Definition / Formula} & \textbf{Compression / streaming role} & \textbf{Notes} \\ \hline
Coding rate constraint & $R_C \leq C_k$ & Prevents packet accumulation and playback stalls & Static $R_C$ causes cliff or leveling effects \\ \hline
MANE & Media-aware middlebox inspecting NAL headers & Drops less important NAL units safely & Useful in managed networks \\ \hline
Push stack & UDP/RTP/RTCP/WebRTC & Ultra-low-latency conversational media & Server-driven \\ \hline
Pull stack & HTTP/DASH/HLS over QUIC/HTTP/3 & Scalable VoD and live streaming & Client-driven \\ \hline
SVC & Base layer plus enhancement layers & One encoded stream can serve many capacities & Poor fit for web-scale CDN caching \\ \hline
DASH & Dynamic Adaptive Streaming over HTTP & Standard HTTP-based ABR architecture & Uses manifest and segment representations \\ \hline
MPD & XML manifest with periods, adaptation sets, representations, segments & Describes available media objects and URLs & Client reads MPD before selecting segments \\ \hline
Closed GOP at segment boundary & Segment starts with IDR & Enables quality switching at boundaries & Open GOPs with cross-segment references are forbidden \\ \hline
CMAF chunks & 200-500 ms sub-segment units & Low-latency DASH/HLS & Parent segment shares one IDR anchor \\ \hline
Throughput & $S(t)=dD(t)/dt$ & Determines download speed at application layer & Lower than physical bitrate $R$ \\ \hline
Average throughput & $\bar{S}=D(T)/T$ & Smooths bursty network measurements & Used by ABR estimation \\ \hline
End-to-end delay & $\sum_k d_{\text{proc}}+d_{\text{queue}}+d_{\text{tx}}+d_{\text{prop}}$ & Determines latency & Propagation and queuing remain structural limits \\ \hline
Playout buffer & $B(t)$ seconds of received but unplayed video & Absorbs jitter and throughput drops & Starvation at $B(t)=0$ causes stalls \\ \hline
Buffer dynamics & $\frac{dB}{dt}=\frac{S}{R_C}-1$ during playout; $\frac{S}{R_C}$ during rebuffering & Fluid model for playback stability & Core ABR system model \\ \hline
Download time & $T_D(n,k)=T_S R_C(k)/S_n$ & Predicts whether segment arrives in time & Assumes average segment throughput $S_n$ \\ \hline
QoE segment metric & $J(n)=\lambda_1 k_n-\lambda_2|k_n-k_{n-1}|-\phi(\Delta_n)$ & Balances quality, switching, and stalls & $\phi$ has large jump for any stall \\ \hline
Sequence QoE & $J=\sum_{n=1}^{N}J(n)-\lambda_3T_{ST}$ & Adds startup delay penalty & $\lambda_i$ chosen from subjective studies \\ \hline
Throughput-based ABR & Select bitrate from throughput estimate & Avoids downloads longer than playout & PANDA, FESTIVE, SQUAD \\ \hline
Buffer-based ABR & Select bitrate from buffer occupancy & Uses buffer as stability indicator & BBA, BOLA \\ \hline
Hybrid ABR & Uses throughput and buffer & Combines prediction with stability & ELASTIC, MPC, ABMA+, DYNAMIC \\ \hline
\end{tabular}
\end{center}

